% Oversells «Fotoalbom» shabloni (XeLaTeX yoki LuaLaTeX).
% \begin{rasmtori}[nisbat]{n} \torrasm{rasm}{izoh} ... \end{rasmtori} — n ta rasm
% qatorda, hammasi bir xil balandlikda (eni × nisbat). Butun betli rasm:
% \rasmramka[joy=tola]{rasm}{}. Koʻrinish sahifa.sty da («Sahifa» tugmasi).
\documentclass[11pt]{article}
\usepackage[paperwidth=240mm,paperheight=170mm]{geometry}
\usepackage{fontspec}
\setmainfont{Noto Sans}
\usepackage{babel}
\babelprovide[import=uz-Latn, main]{uzbek}
\def\uzbekhyphenmins{22}
\usepackage{sahifa} % Oversells sahifa dizayni (sahifa.sty)
\raggedbottom

\begin{document}

\section*{Tabiat}
\begin{uslub}{kirish}
Bu albomga yozgi safarlarimizdan eng yaxshi suratlarni yigʻdik. Har bir
rasm ostida — qayerda va qachon olingani.
\end{uslub}

\begin{rasmtori}{3}
\torrasm{rasmlar/tog.png}{Chimyon togʻlari, iyul}
\torrasm{rasmlar/bog.png}{Bogʻdagi ertalab}
\torrasm{rasmlar/kol.png}{Koʻl boʻyida}
\end{rasmtori}

\begin{rasmtori}[0.45]{2}
\torrasm{rasmlar/qishloq.png}{Buvimning uyi}
\torrasm{rasmlar/shom.png}{Shom payti}
\end{rasmtori}

\clearpage
\section*{Dala}
\begin{rasmtori}{4}
\torrasm{rasmlar/dala.png}{Bugʻdoyzor}
\torrasm{rasmlar/tog.png}{Choʻqqi}
\torrasm{rasmlar/kol.png}{Suv}
\torrasm{rasmlar/bog.png}{Daraxtlar}
\torrasm{rasmlar/shom.png}{Qushlar}
\torrasm{rasmlar/qishloq.png}{Qishloq}
\torrasm{rasmlar/dala.png}{Yoʻl}
\torrasm{rasmlar/bog.png}{Soya}
\end{rasmtori}

\rasmramka[joy=tola]{rasmlar/dala.png}{}

\end{document}
